\textbf{Magneticky měkké materiály} mají malou koercivitu a úzkou hysterezní smyčku:
\begin{itemize}
\item technicky čisté Fe -- stejnosměrné magnetické obvody;
\item Fe--Si plechy -- transformátory a elektrické stroje;
\item Fe--Ni slitiny (permalloy, supermalloy, mumetal) -- jádra transformátorů a tlumivek, magnetické stínění;
\item sendust (Fe--Al--Si) a měkké ferity -- magnetická jádra cívek a tlumivek.
\end{itemize}
\textbf{Magneticky tvrdé materiály} mají velkou koercivitu a širokou smyčku. Patří sem oceli legované W, Cr či Co, slitiny AlNi, AlNiCo, CuNiFe a CuNiCo; používají se pro \textbf{permanentní magnety}.
